\section{Introduction}

\label{sec:introduction}

The conversion of natural habitats to human-dominated land uses is the single largest driver of biodiversity decline, responsible for more species extinctions than climate change, pollution, and overexploitation combined \citep{maxwell2016biodiversity}.
Tropical deforestation alone eliminates an estimated 10 million hectares of forest annually \citep{fao2020global}, while less conspicuous forms of habitat degradation---selective logging, agricultural intensification, infrastructure development---erode ecosystem integrity across all biomes \citep{watson2018decline}.

Satellite remote sensing provides the only feasible means of monitoring land use change at scales relevant to global biodiversity conservation.
The Landsat archive offers 50 years of multispectral imagery at 30-meter resolution \citep{wulder2019current}.
The Sentinel-2 constellation provides 10-meter resolution with 5-day revisit frequency \citep{drusch2012sentinel}.
Commercial constellations (Planet, Maxar) achieve sub-meter resolution with daily coverage.

Despite abundant data, three challenges limit the effectiveness of satellite-based ecological monitoring.

First, \emph{change detection accuracy}.
Traditional approaches based on vegetation indices (NDVI differencing) or pixel-level classification comparison produce high false-positive rates due to phenological variation, atmospheric artifacts, and cloud contamination \citep{zhu2017change}.
Deep learning methods have improved accuracy but require extensive labeled training data that is expensive to produce for diverse global ecosystems \citep{yuan2021deep}.

Second, \emph{ecological interpretation}.
Raw change detection identifies where land cover has changed but not what the change means for biodiversity.
A detected change could represent catastrophic deforestation, natural seasonal variation, prescribed burning, or sustainable agriculture---outcomes with vastly different conservation implications \citep{redford2003empty}.
Translating remote sensing signals into ecologically meaningful assessments requires integration with landscape ecology principles.

Third, \emph{actionable timeliness}.
Global forest monitoring systems such as Global Forest Watch \citep{hansen2013high} provide monthly to annual updates.
For conservation enforcement, this latency is often too slow: illegal deforestation operations can clear hundreds of hectares within days, and by the time satellite-based alerts arrive, irreversible damage has occurred \citep{finer2018combating}.

We present ZooSat, a deep learning pipeline for continuous ecological monitoring that addresses all three challenges.
Our contributions are:

\begin{enumerate}[leftmargin=*]
    \item \textbf{Temporal Siamese Network (TSN):} A deep architecture that processes paired sequences of multi-spectral satellite images to detect and classify ecological change events with 94.2\% accuracy, handling cloud contamination and phenological variation through temporal attention mechanisms (\Cref{sec:tsn}).
    \item \textbf{Habitat Fragmentation Index (HFI):} A landscape-scale metric computed from detected changes that quantifies habitat connectivity degradation in real-time, providing a leading indicator of biodiversity decline (\Cref{sec:fragmentation}).
    \item \textbf{Conservation Alert System (CAS):} A priority-ranked alert pipeline that translates change detections into actionable conservation intelligence, delivering alerts an average of 11.2 days faster than existing systems (\Cref{sec:alerts}).
\end{enumerate}

% ============================================================
% 2. Related Work
% ============================================================
